\documentclass[reqno,10pt,a4paper]{amsart}
\usepackage{mathtools,amssymb,microtype}
\usepackage[T1]{fontenc}
\usepackage{lmodern}
\usepackage[a4paper,textwidth=15.5cm,hcentering,top=1.8cm,bottom=2.1cm,headheight=10pt,headsep=14pt]{geometry}
\usepackage{xcolor}
\definecolor{linkblue}{RGB}{0,61,105}
\usepackage[colorlinks=true,linkcolor=linkblue,citecolor=linkblue,urlcolor=linkblue]{hyperref}
\setlength{\emergencystretch}{1em}
\numberwithin{equation}{section}
\newcommand{\R}{\mathbb R}
\newcommand{\T}{\mathbb T}
\linespread{1.15}
\title[Long-time behavior and stability]{Long-time behavior and stability of\\ kinetic and dispersive equations}
\author{Wenrui Huang}
\date{}
\hypersetup{pdftitle={Project Summary -- Long-time behavior and stability of kinetic and dispersive equations},pdfauthor={Wenrui Huang}}
\begin{document}
\maketitle
\begin{center}
{\small Herchel Smith Research Fellowship in Pure Mathematics \enspace$\cdot$\enspace Project Summary}
\end{center}
\smallskip
\noindent\textbf{Research areas.} Nonlinear partial differential equations, kinetic theory, and dispersive equations.
\medskip

My research concerns the long-time behavior of kinetic and dispersive equations. Once global existence of solutions is understood, two basic dynamical questions arise. How do solutions behave at large times, and are stationary solutions stable under small perturbations? I study these questions for collisionless plasmas and for models in which particles interact with a dispersive wave.

\subsection*{Scattering for particles and waves.}
In joint work with Hyunwoo Kwon, I studied how the range of the interaction changes the trajectories in the Vlasov--Riesz equation. With Mengyi Xie, I studied the corresponding problem for the Hartree equation and for a coupled Vlasov--Hartree system. In the coupled system, the particle density generates the potential acting on the wave, and the wave density generates the force on the particles. For slowly decaying interactions, we found that a common nonlinear action determines both the correction to the particle trajectories and the phase correction to the wave.

An external potential changes the underlying motion even in the absence of interactions. For Vlasov--Poisson with a repulsive harmonic potential, Kwon and I introduced a hyperbolic lens transform and constructed wave operators. In one dimension, my subsequent work gives modified scattering even for large data when the self-interaction is repulsive. I have also constructed modified wave operators for the Hartree equation with a repulsive Coulomb potential. I would like to study the corresponding Vlasov dynamics in an \emph{attractive} harmonic potential, where our lens transform still applies. A broader goal is to identify the radial external potentials for which solutions exhibit scattering or modified scattering, beginning with radially symmetric solutions. At fixed angular momentum, the radial motion under the external Hamiltonian is integrable. I would like to understand the correction to this motion due to the evolving density.

\subsection*{Plasma--wall interactions.}
With Benoit Pausader and Masahiro Suzuki, I studied a dilute plasma in an unbounded convex domain with a perfectly conducting, absorbing wall. We identified the cone of velocities compatible with escape to infinity. Its geometry and the limiting particle distribution together determine the leading electric field along escaping trajectories. I propose to study a perfectly reflecting wall, where particles undergo specular reflection and can encounter the boundary repeatedly. The goal is to understand how these encounters affect global regularity, asymptotic stability, and the long-time distribution of particles.

\subsection*{Stability of small BGK waves.}
Bernstein--Greene--Kruskal (BGK) waves are spatially inhomogeneous stationary states of the Vlasov--Poisson system. With Dongfen Bian, Emmanuel Grenier, and Benoit Pausader, I obtained a criterion for the linear stability or instability of small one-well waves bifurcating near a generic marginal homogeneous equilibrium. On the stable branch, the squared spatial norm of the dynamical electric field has a finite time integral. I am working toward quantitative damping estimates and, under suitable regularity assumptions, exponential decay of this field. A longer-term goal is to establish nonlinear stability. On the unstable branch, the next question is whether the growing linear mode leads to nonlinear instability.

At Cambridge, I would welcome the opportunity to discuss kinetic theorey with Cl\'ement Mouhot, and nonlinear dispersive equations with Pierre Rapha\"el.

\end{document}
